\documentclass[11pt,a4paper]{report}

\usepackage[utf8]{inputenc}
\usepackage[T1]{fontenc}
\usepackage[danish]{babel}
\usepackage[a4paper,margin=2.5cm]{geometry}
\usepackage{amsmath}
\usepackage{graphicx}
\usepackage{float}
\usepackage{booktabs}
\usepackage{listings}
\usepackage{xcolor}
\usepackage{hyperref}
\usepackage{caption}

\definecolor{codebg}{RGB}{245,245,245}
\lstset{
    backgroundcolor=\color{codebg},
    basicstyle=\ttfamily\small,
    breaklines=true,
    frame=single,
    numbers=left,
    numberstyle=\tiny,
    keywordstyle=\color{blue},
    commentstyle=\color{gray},
    showstringspaces=false,
    columns=fullflexible
}

\setlength{\parindent}{0pt}
\setlength{\parskip}{0.5em}

\title{\vspace{1cm}\Huge Miniprojekt A\\[0.3cm]\Large Sampling og analyse af digitale signaler}
\author{\textbf{Navn:} \\ \textbf{AU-ID:} \\ \textbf{Afleveringsdato:} \today}
\date{}

\begin{document}

\maketitle

\section{Opgave 1: Find antal samples N for alle tre signaler}
\subsection{Opgavebeskrivelse}
Beskriv kort, hvad opgaven går ud på.

\subsection{Forberedelse}
Skriv dine forberedende overvejelser her. Hvad skal du måle, og hvordan vil du finde det?

\subsection{Kode}
\begin{lstlisting}[language=Python]
# Indsæt koden fra Python-notebooket til denne opgave her
# Eksempel:
# N_s1 = len(s1)
# N_s2 = len(s2)
# N_s3 = len(s3)
\end{lstlisting}

\subsection{Resultat}
Skriv resultaterne her. Indsæt værdier, grafer eller en kort summary af det, du fandt.

\subsection{Diskussion og konklusion}
Skriv en kort diskussion af resultatet og konklusionen for opgaven.

\newpage

\section{Opgave 2: Plot alle tre signaler med korrekte tidsakser}
\subsection{Opgavebeskrivelse}
Beskriv kort, hvad opgaven går ud på.

\subsection{Forberedelse}
Skriv dine forberedende overvejelser her.

\subsection{Kode}
\begin{lstlisting}[language=Python]
# Indsæt koden fra Python-notebooket til denne opgave her
# Eksempel:
# t1 = np.arange(len(s1)) / fs
# plt.plot(t1, s1)
\end{lstlisting}

\subsection{Resultat}
Skriv resultaterne her.

\subsection{Diskussion og konklusion}
Skriv en kort diskussion af resultatet og konklusionen for opgaven.

\newpage

\section{Opgave 3: Find max, min, gennemsnit, RMS og effekt for alle tre signaler}
\subsection{Opgavebeskrivelse}
Beskriv kort, hvad opgaven går ud på.

\subsection{Forberedelse}
Skriv dine forberedende overvejelser her.

\subsection{Kode}
\begin{lstlisting}[language=Python]
# Indsæt koden fra Python-notebooket til denne opgave her
# Eksempel:
# maks = np.max(s1)
# rms = np.sqrt(np.mean(s1**2))
\end{lstlisting}

\subsection{Resultat}
Skriv resultaterne her.

\subsection{Diskussion og konklusion}
Skriv en kort diskussion af resultatet og konklusionen for opgaven.

\newpage

\section{Opgave 4: Find crest-faktorerne for de tre signaler og sammenlign}
\subsection{Opgavebeskrivelse}
Beskriv kort, hvad opgaven går ud på.

\subsection{Forberedelse}
Skriv dine forberedende overvejelser her.

\subsection{Kode}
\begin{lstlisting}[language=Python]
# Indsæt koden fra Python-notebooket til denne opgave her
# Eksempel:
# cf_s1 = np.max(np.abs(s1)) / np.sqrt(np.mean(s1**2))
\end{lstlisting}

\subsection{Resultat}
Skriv resultaterne her.

\subsection{Diskussion og konklusion}
Skriv en kort diskussion af resultatet og konklusionen for opgaven.

\newpage

\section{Opgave 5: Nedsampling af s1 med faktor 4}
\subsection{Opgavebeskrivelse}
Beskriv kort, hvad opgaven går ud på.

\subsection{Forberedelse}
Skriv dine forberedende overvejelser her.

\subsection{Kode}
\begin{lstlisting}[language=Python]
# Indsæt koden fra Python-notebooket til denne opgave her
# Eksempel:
# s1_down = s1[::4]
\end{lstlisting}

\subsection{Resultat}
Skriv resultaterne her.

\subsection{Diskussion og konklusion}
Skriv en kort diskussion af resultatet og konklusionen for opgaven.

\newpage

\section{Opgave 6: Lytning og beskrivelse af forskellen mellem originalen og den nedsamplede version}
\subsection{Opgavebeskrivelse}
Beskriv kort, hvad opgaven går ud på.

\subsection{Forberedelse}
Skriv dine forberedende overvejelser her.

\subsection{Kode}
\begin{lstlisting}[language=Python]
# Indsæt koden fra Python-notebooket til denne opgave her
# Eksempel:
# Audio(s1, rate=fs)
# Audio(s1_down, rate=fs/4)
\end{lstlisting}

\subsection{Resultat}
Skriv resultaterne her.

\subsection{Diskussion og konklusion}
Skriv en kort diskussion af resultatet og konklusionen for opgaven.

\newpage

\section{Opgave 7: Sammenlign de to stereo-kanaler i s2}
\subsection{Opgavebeskrivelse}
Beskriv kort, hvad opgaven går ud på.

\subsection{Forberedelse}
Skriv dine forberedende overvejelser her.

\subsection{Kode}
\begin{lstlisting}[language=Python]
# Indsæt koden fra Python-notebooket til denne opgave her
# Eksempel:
# left = s2[:, 0]
# right = s2[:, 1]
# plt.plot(left)
# plt.plot(right)
\end{lstlisting}

\subsection{Resultat}
Skriv resultaterne her.

\subsection{Diskussion og konklusion}
Skriv en kort diskussion af resultatet og konklusionen for opgaven.

\newpage

\section{Opgave 8: Kvantisering af s2 med 4 bit}
\subsection{Opgavebeskrivelse}
Beskriv kort, hvad opgaven går ud på.

\subsection{Forberedelse}
Skriv dine forberedende overvejelser her.

\subsection{Kode}
\begin{lstlisting}[language=Python]
# Indsæt koden fra Python-notebooket til denne opgave her
# Eksempel:
# levels = 2**4
# s2_q = np.round(s2_left * (levels - 1)) / (levels - 1)
\end{lstlisting}

\subsection{Resultat}
Skriv resultaterne her.

\subsection{Diskussion og konklusion}
Skriv en kort diskussion af resultatet og konklusionen for opgaven.

\newpage

\section{Opgave 9: Fade-out på s3 med eksponentielt aftagende envelope}
\subsection{Opgavebeskrivelse}
Beskriv kort, hvad opgaven går ud på.

\subsection{Forberedelse}
Skriv dine forberedende overvejelser her.

\subsection{Kode}
\begin{lstlisting}[language=Python]
# Indsæt koden fra Python-notebooket til denne opgave her
# Eksempel:
# N = len(s3)
# start = int(2 * N / 3)
# envelope[start:] = np.exp(-alpha * np.arange(N - start))
\end{lstlisting}

\subsection{Resultat}
Skriv resultaterne her.

\subsection{Diskussion og konklusion}
Skriv en kort diskussion af resultatet og konklusionen for opgaven.

\newpage

\section{Opgave 10: Lav et enkelt ekko med delay på 100 ms}
\subsection{Opgavebeskrivelse}
Beskriv kort, hvad opgaven går ud på.

\subsection{Forberedelse}
Skriv dine forberedende overvejelser her.

\subsection{Kode}
\begin{lstlisting}[language=Python]
# Indsæt koden fra Python-notebooket til denne opgave her
# Eksempel:
# delay = int(0.1 * fs)
# y = x.copy()
# y[delay:] += 0.5 * x[:-delay]
\end{lstlisting}

\subsection{Resultat}
Skriv resultaterne her.

\subsection{Diskussion og konklusion}
Skriv en kort diskussion af resultatet og konklusionen for opgaven.

\end{document}